\documentclass[11pt]{article}
\usepackage[review]{acl}
\usepackage{times,latexsym}
\usepackage[T1]{fontenc}
\usepackage[utf8]{inputenc}
\usepackage{microtype}
\usepackage{graphicx,booktabs,amsmath,amssymb,tabularx}
\usepackage{url}
\hypersetup{hidelinks,pdftitle={No Detections, Visible Nudity},pdfauthor={Anonymous research draft}}

\title{No Detections, Visible Nudity:\\Evidence-First Review of an Original Image}
\author{Anonymous authors}
\newcommand{\method}{Evidence-First Audit}
\newcommand{\yes}{\textsf{yes}}
\newcommand{\no}{\textsf{no}}

\begin{document}
\maketitle

\begin{abstract}
An original image contains an unclothed figure in an artwork displayed on an
easel, yet NudeNet returns no detections and two public NSFW classifiers assign
scores of 0.000199 and 0.003657. This observation requires neither background
replacement nor adversarial image modification: it occurs on the supplied image
itself. We study the resulting audit blind spot: an empty detector output or a
low global safety score cannot establish that visible nudity is absent.
We propose an evidence-first review protocol that describes visible content
before applying a contextual policy judgment. On the identical original input,
GPT-5.6 Luna identifies the unclothed depiction while separately reporting no
clear explicit anatomical detail and no sexual activity. In an exploratory
matched set of 18 inputs, a structured review agrees with all reference content labels. This single-source case study motivates
an auditable separation of detection, visibility, and policy, rather than a
claim that artistic nudity is necessarily prohibited or that a new classifier
has achieved general superiority.
\end{abstract}

\begin{figure*}[t]
\centering
\begin{minipage}[c]{0.34\linewidth}
\centering
\includegraphics[width=\linewidth]{figures/original_input.jpg}
\end{minipage}\hfill
\begin{minipage}[c]{0.62\linewidth}
\small
\textbf{The supplied original, not a neutral-background substitute.}

\begin{tabularx}{\linewidth}{@{}lX@{}}
\toprule
System & Actual output on this image \\
\midrule
NudeNet 3.4.2 & \texttt{detections = []} \\
Falconsai & NSFW score: 0.00019937 \\
AdamCodd & NSFW score: 0.00365668 \\
\midrule
Luna, structured & Unclothed depiction: \yes \\
review & Explicit detail: \no; activity: \no \\
\bottomrule
\end{tabularx}

\textbf{Central finding.} The detector signals do not surface the visible
unclothed depiction. A direct, content-specific review of the same image does.
\end{minipage}
\caption{\textbf{The original-image finding that motivates the paper.}
The displayed image is a byte-identical copy of the supplied input, not an
edited or re-composited variant. Each system uses its native preprocessing.
NSFW scores, body-part detections, and content attributes have different
semantics; this is an evidence comparison, not a common-label accuracy table.}
\label{fig:original}
\end{figure*}

\section{Introduction: the missed content is already in the original}
\label{sec:intro}

\textbf{The starting point is one concrete image, not a transformation.}
Figure~\ref{fig:original} shows the original artwork-on-easel input. Its central
artwork visibly depicts an unclothed human figure. Direct evaluation of the
unmodified file produces an empty NudeNet result and very low NSFW scores from
two public image classifiers. A reviewer who treats those outputs as evidence
that the image contains no nudity would miss content that is visibly present.
This is the central problem studied here.

Safety systems often compress different questions into a single signal:
is a person depicted without clothing; are explicit anatomical details visible;
is there sexual activity; and should a particular service allow the image?
These questions are not interchangeable. A stylized artwork can contain an
unclothed depiction without depicting sexual activity, and an applicable policy
may permit it. Conversely, permission does not erase what is depicted. A useful
audit therefore needs a content record that survives the final allow/reject
decision.

Our proposed response is \emph{not} to manipulate the image until a detector
fires. \method{} asks a vision-language reviewer to inspect the original image
and return separate content attributes, visual evidence, and uncertainty.
Policy assessment is a subsequent operation. We instantiate this protocol
using the requested GPT-5.6 Luna model and evaluate it on the same original
file used by the detectors. It recognizes the unclothed depiction without
conflating it with sexual activity. No crop, neutral canvas, or replacement
background is needed for this primary result.

\paragraph{Contributions and scope.}
We contribute (i) a reproducible original-image case with raw detector outputs
and byte-level input identity; (ii) an executable evidence-first audit protocol
that separates visible content from policy; and (iii) a matched exploratory
evaluation of structured content recognition. Additional background experiments are diagnostic
supplementary evidence, not the premise of the paper. Our evidence supports a
case-level audit blind spot, not a population-wide failure rate, a newly trained
model, or a demonstrated violation of a closed-source service's policy.

\section{Original-image evidence}
\label{sec:case}

\subsection{Input identity and reference content}

The supplied file is a JPEG image of size $848\times1264$, despite the
\texttt{.png} suffix in its original pathname. Its SHA-256 is recorded in
Appendix~\ref{app:repro}. The local asset, all three direct detector records,
and the anonymous Luna input \texttt{S04} have this same hash. Consequently,
differences below cannot be attributed to selecting a neutral-background
variant or substituting a different image for the vision-language review.

The case-level content reference is \emph{an unclothed human depiction within
the artwork}. No visible sexual activity is annotated. This reference was
assigned by the research agent from the visible image, not by an independent
human panel. The Luna review separately reports the presence of the
unclothed depiction and the absence of clear explicit anatomical detail and
sexual activity. We do not infer a real person's identity or the provenance
of the depicted scene.

\subsection{What the existing detectors actually return}

We run NudeNet 3.4.2 with its 320n ONNX model and native preprocessing.
The raw result is an empty list: no body-part detections are returned.
The implementation uses a raw candidate threshold of 0.20, an NMS score
threshold of 0.25, and NMS IoU threshold of 0.45. We preserve the empty output
as such; we do not manufacture a body-part confidence or interpret zero
detected box area as zero visible body area.

We also run a community FP32 ONNX export of
\texttt{Falconsai/nsfw\_image\_detection} and the FP32 ONNX model from
\texttt{AdamCodd/vit-base-nsfw-detector}. They assign NSFW scores
$0.0001993742$ and $0.0036566833$, respectively. Neither raises an NSFW
alarm at threshold 0.50, or at any of the recorded thresholds
$\{0.20,0.35,0.50,0.65,0.80\}$. These are model class scores, not calibrated
probabilities that the image lacks nudity. Model revisions and preprocessing
are retained with the raw outputs.

\paragraph{The precise failure claim.}
The tested detector signals \emph{do not surface the unclothed depiction in
this original image}. If these signals are used as the sole check for whether
such content exists, the resulting audit has a false negative with respect to
that content task. This does not by itself establish an error under either
classifier's native NSFW policy, nor a body-part localization error against
independently annotated NudeNet boxes. The distinction makes the warning more
actionable: do not use a negative safety signal as a complete content inventory.

\section{Method: describe the evidence before judging the context}
\label{sec:method}

\subsection{A content record, not another undifferentiated ``sex'' score}

For an original image $x$, the first stage produces
\begin{equation}
e(x)=(u,d,a,m,r,q), \qquad u,d,a\in\{0,1,\bot\},
\label{eq:record}
\end{equation}
where $u$ denotes an unclothed human depiction, $d$ clear explicit anatomical
detail, $a$ sexual activity, $m$ the depicted medium, $r$ a brief visual-evidence
description, and $q$ uncertainty. The symbol $\bot$ means the content cannot be
determined from the visible image. It must not silently become a negative
label. We retain these attributes even when the eventual policy permits the
image.

The second stage applies the actual deployment policy $\pi$:
\begin{equation}
z=\pi\big(e(x),c(x)\big),
\label{eq:policy}
\end{equation}
where $c(x)$ is relevant context and $z$ is an allow, review, or reject outcome.
The first stage was executed in this study; the second is a proposed integration
contract, not an evaluated production policy engine. The design does not
require artistic or medical context to be ignored when making a policy
decision. It requires that such context not replace an account of the visible
content.

\begin{table*}[tp]
\centering
\small
\begin{tabularx}{\linewidth}{@{}lXX@{}}
\toprule
Stage & Input and output & Required behavior \\
\midrule
Evidence review & Original image $x\rightarrow e(x)$ & Inspect depicted content,
including artwork within the scene; separate the three attributes. \\
Policy assessment & $e(x)$, context, policy $\rightarrow z$ & Allow permitted art
when appropriate; do not erase a positive content attribute. \\
Audit logging & Input hash, $e(x)$, $z$, versions & Make disagreement inspectable;
retain abstentions and the actual detector output. \\
\bottomrule
\end{tabularx}
\caption{\method{} as a review protocol. Only the evidence-review stage has
been empirically evaluated here. No automatic crop or image transformation is
part of the primary original-image experiment.}
\label{tab:method}
\end{table*}

\subsection{Executed structured review}

An agent context was explicitly assigned \texttt{gpt-5.6-luna} and
inspected 18 anonymous image files in a fixed order using the image-viewing
tool. Instructions forbade access to reference labels, baseline scores, and
the manuscript. This was instruction-level label blinding, not operating-system
isolation. The original file appeared as \texttt{S04}.

The instruction asked for three separate content attributes:
\begin{quote}\small
Based only on what is visible in each image, determine whether it contains
an unclothed person or depiction, clear explicit anatomical detail, and sexual
activity. Do not merge these three judgments.
\end{quote}
The executed instructions were in Chinese; the wording above is an English
translation. The original protocol and responses are retained in the research
files. This structured review implements content decomposition; we do not
present an ambiguous one-word ``sex?'' query as an evaluated baseline.

\section{Evaluation: does the method see the original image?}
\label{sec:results}

\subsection{Primary result: direct review recovers the visible content}

\begin{table}[ht]
\centering
\small
\begin{tabularx}{\linewidth}{@{}Xccc@{}}
\toprule
Review & Unclothed & Detail & Activity \\
\midrule
Luna, structured & \yes & \no & \no \\
\bottomrule
\end{tabularx}
\caption{\textbf{Same original input, no image modification.} The review
identifies the content not surfaced by the detector signals in
Figure~\ref{fig:original}. A single image provides a case outcome, not an
accuracy estimate.}
\label{tab:original-review}
\end{table}

The structured review describes an unclothed figure in the woodblock-style
artwork on the easel and reports no uncertainty. This is an actual content
judgment on \texttt{S04}, not a prediction inferred from scores on a crop or
a background-replaced image.

The operational improvement is an explicit content record: the system can
acknowledge that nudity is depicted while retaining separate negative findings
for clear explicit detail and sexual activity. It need not choose between
``everything is safe, therefore nothing is depicted'' and ``a nude depiction
exists, therefore reject.'' The original case demonstrates this capability;
it does not isolate which internal model mechanism enables it.

\subsection{Exploratory matched comparison}

To check that the reviewer did not simply label every input as containing an
unclothed figure, we include five benign sources: coffee, a cat, a rocket, a
clothed astronaut, and a clothed camera operator. The exploratory set comprises
the original, 12 background/scale variants of its artwork, and one condition
from each benign source. This yields 13 positive conditions and five negatives,
but only one positive source and five negative sources.

\begin{table*}[tp]
\centering
\small
\begin{tabular}{@{}lrrrr@{}}
\toprule
Method, same 18 inputs & TP/FN & FP/TN & Accuracy & Recall \\
\midrule
CLIP, generic full-image reference & 6/7 & 2/3 & 50.00\% & 46.15\% \\
CLIP, art-aware full-image reference & 10/3 & 0/5 & 83.33\% & 76.92\% \\
Luna, basic structured review & 13/0 & 0/5 & 100.00\% & 100.00\% \\
\bottomrule
\end{tabular}
\caption{Descriptive recognition results for \emph{unclothed human depictions},
not pornography or policy violations. The inputs and reference labels are
identical across rows. The correlated variants are not independent examples.
CLIP is an auxiliary task-matched reference, not the paper's proposed method.}
\label{tab:matched}
\end{table*}

The Luna review agrees with all 18 reference labels and correctly leaves the
five benign conditions negative. Its false-positive rate is zero on these
five conditions.

These 18/18 results must not be read as general 100\% accuracy. The positive
conditions derive from the same artwork, reference labels come from the
research agent, and the review has only one sequential evaluation
context. There are no sexual-activity positives, no medical images, and no
independent human adjudication. We therefore cannot estimate sexual-activity
recall, medical-domain accuracy, or source-level generalization.

\section{Measurement beyond attack success rate}
\label{sec:metrics}

The original-image question is not an attack-success-rate (ASR) question.
ASR requires a defined target policy and an attack outcome. It neither measures
how much content is visible nor distinguishes an undetected depiction from
an allowed artwork. We instead report the following complementary quantities.

\paragraph{Content recognition.}
For an explicit content reference, report recall, false-positive rate (FPR),
accuracy, and balanced accuracy (BA), together with raw confusion counts:
\begin{equation}
\begin{aligned}
\mathrm{Recall}&=\frac{TP}{TP+FN},\quad
\mathrm{FPR}=\frac{FP}{FP+TN},\\
\mathrm{BA}&=\frac{\mathrm{Recall}+1-\mathrm{FPR}}{2}.
\end{aligned}
\end{equation}
If an attribute has no positive references, its recall is unidentifiable, not
zero. Abstentions should be reported separately with coverage; none occurred
in the recorded Luna review. Native NSFW scores are not inserted into
this content-accuracy table as though the tasks were identical.

\paragraph{Visibility and input resolution.}
For a research annotation $M$ and preprocessing transform $T_j$ of model $j$,
we record a frame-area fraction and effective reference pixels:
\begin{equation}
A(x)=\frac{|M|}{WH},\qquad E_j(x)=|T_j(M)|.
\end{equation}
The annotation is a coarse visible torso-and-limb region, not a segmentation
of all skin or a ``pornography area.'' These measures describe the evidence
available to a model; they do not define sexual severity. The supplementary
scale controls demonstrate why identical semantic labels may occupy very
different numbers of model-input pixels.

\paragraph{Continuous signals and policy decisions.}
Keep the actual class scores and detections, rather than only thresholded
outcomes. A score shift is not necessarily a content-label change, and neither
necessarily changes the policy outcome. For future ordinal severity judgments,
an independently annotated rubric and agreement study would be needed before
reporting ordinal error or rank correlation. We do not invent such a severity
scale from this one image.

\section{Related work}

Safe Latent Diffusion and I2P evaluate inappropriate generations using explicit
detector-based criteria \citep{sld}. NudeNet provides body-part detections and
scores \citep{nudenet}, whereas SafeSearch exposes several content likelihood
categories \citep{safesearch}. These outputs are useful but are not direct
measurements of visible surface area or interchangeable definitions of nudity.
UnsafeBench examines safety classifiers on real-world and generated images,
motivating attention to evaluation domain \citep{unsafebench}.

Artistic nudity has already been studied explicitly in content moderation
\citep{artcentric}. Our work does not claim to discover that art is a difficult
domain, nor to introduce the first use of CLIP for artistic content. Likewise,
LlavaGuard establishes prior work on structured vision-language safety
assessment \citep{llavaguard}. Our narrower contribution is an original-image
case with verified input identity, direct evidence-first review, and matched content-recognition diagnostics. Six-CD studies
concept removal and preservation \citep{sixcd}; content preservation and
visible nudity measurement remain distinct tasks. We discuss these works as
related literature, not as models evaluated in our experiment.

\section{Discussion and conclusion}

\paragraph{What caused the empty result?}
Stylization, small effective visual regions, domain mismatch, background
features, and label definitions are plausible explanations. The direct-image
experiment does not identify their causal contributions. Supplementary
background and crop controls reveal score sensitivity and false-positive
trade-offs, but cannot establish an internal artistic-exemption mechanism.
The paper's central observation does not depend on such an attribution.

\paragraph{Practical implication.}
If a product needs to know whether nudity is depicted, an empty detector result
should not close the content audit. Preserve evidence fields alongside policy
decisions, and assess review mechanisms on both missed depictions and benign
artwork or object false positives. Any escalation policy should be evaluated
under an explicit review budget rather than assumed to work because one
example is recovered. The source image remains local because public
redistribution rights have not been established; no new explicit content was
generated for this study.

\paragraph{Conclusion.}
\textbf{The supplied original already contains the finding: visible artistic
nudity, no NudeNet detections, and very low NSFW scores.} Direct structured
review of that same image records the unclothed depiction without mislabeling
it as sexual activity. This motivates evidence-first auditing, while the single-source design delimits what has actually been demonstrated. Background transformations are supplementary diagnostics,
not the central claim.

\section*{Limitations}
\paragraph{What is needed for a stronger method claim?}
A source-disjoint dataset, independently adjudicated content labels, repeated
reviews, task-matched baselines, and cost/latency measurements are required.
NudeNet should additionally be evaluated against its native body-part labels.
The Luna identifier records the model requested through the agent interface;
we lack an independently verifiable serving snapshot and calibrated
probabilities. There is no new training, automatic artwork localization,
deployed policy engine, or demonstrated end-to-end superiority in this study.

\section*{Reproducibility statement}
The local research bundle contains the original input, its hash, frozen model
metadata, raw detector outputs, the structured Luna review file, sample mappings, and
analysis code. The accompanying verification script checks all 211 image hashes
against detector records and all 18 review-input mappings before validating the
paper's key figures. No missing experiment is converted into a zero score.

\bibliography{references}

\appendix
\section{Supplementary diagnostics: background is not the primary input}
\label{app:background}

These experiments explore the original finding; they do not replace the
original-image evaluation. We extract the complete artwork using the manually
specified rectangle $[120,350,744,924]$, with right and bottom endpoints
excluded. It includes both the figure and the artwork's wave pattern. This is
an oracle region, not automatically discovered anatomy.

We use one neutral background and three natural scenes (living room, office,
and bookshelves). Each natural scene additionally has a mean-color and an
$8\times8$ shuffled-block control, producing ten background versions. Six
foreground sources and three scales yield $6\times10\times3=180$ composites:
30 artwork conditions and 150 benign conditions. Eighteen auxiliary crops,
ten background-only images, the original, and two historical outputs bring
the checked inventory to 211 images. All 180 composites retain identical
foreground pixels across backgrounds at the same scale. None is represented
as a newly successful generative edit.

\subsection{Actual detector scores for the 12 natural/reference conditions}

\begin{table*}[tp]
\centering
\small
\begin{tabular}{@{}llrrr@{}}
\toprule
Scale & Background & Falconsai NSFW & AdamCodd NSFW & CLIP art margin \\
\midrule
0.50 & Neutral gray & 0.00064829 & 0.02091475 & -0.00723764 \\
0.50 & Living room & 0.00095839 & 0.00849401 & +0.00565173 \\
0.50 & Office & 0.00016494 & 0.01989171 & +0.00362690 \\
0.50 & Bookshelves & 0.00025379 & 0.03958590 & +0.00208649 \\
0.75 & Neutral gray & 0.00049196 & 0.02898531 & +0.02188858 \\
0.75 & Living room & 0.00294859 & 0.00498540 & +0.00118987 \\
0.75 & Office & 0.00028360 & 0.00910379 & -0.00167774 \\
0.75 & Bookshelves & 0.00019994 & 0.03295755 & +0.00767250 \\
1.00 & Neutral gray & 0.00035368 & 0.03368190 & +0.00596030 \\
1.00 & Living room & 0.00032044 & 0.00437642 & +0.00970811 \\
1.00 & Office & 0.00028989 & 0.01294989 & +0.00945719 \\
1.00 & Bookshelves & 0.00019713 & 0.02003348 & +0.01551880 \\
\bottomrule
\end{tabular}
\caption{Actual supplementary scores. These are altered-background conditions,
\emph{not the original input}. None of these rows is substituted for the
original-image scores in Figure~\ref{fig:original}.}
\end{table*}

CLIP ViT-B/32 uses three classes (unclothed depiction, clothed person, and no
person), each with three fixed descriptions. Its margin is mean cosine
similarity to the unclothed class minus the larger of the other two class
means; the fixed decision is margin $>0$. Generic and art-aware descriptions
are recorded in the frozen specification. The ordinary oracle-crop condition
was added after the initial smoke test and is identified as a diagnostic
rather than a pre-registered primary condition.

At scale 0.75, the same foreground changes from a positive art-aware margin on
gray to a negative one in the office. This is background-associated
sensitivity, not a demonstration of a universal best evasion background.
At scale 0.50 the direction differs. Mean-color and block-shuffle controls also
show that low-level statistics can matter; the experiment does not isolate a
pure semantic-context effect.

\subsection{Recovery has false-positive costs}

\begin{table*}[tp]
\centering
\small
\begin{tabular}{@{}lrrrrr@{}}
\toprule
CLIP condition, 180 composites & TP/FN & FP/TN & Recall & FPR & Accuracy \\
\midrule
Generic full image & 16/14 & 15/135 & 53.33\% & 10.00\% & 83.89\% \\
Art-aware full image & 20/10 & 0/150 & 66.67\% & 0.00\% & 94.44\% \\
Generic oracle artwork & 30/0 & 90/60 & 100.00\% & 60.00\% & 50.00\% \\
Art-aware oracle artwork & 30/0 & 30/120 & 100.00\% & 20.00\% & 83.33\% \\
\bottomrule
\end{tabular}
\caption{The complete supplementary crop diagnostic, including failures.
The artwork region is manually known. Conditions are correlated within source.}
\end{table*}

All 30 false positives of the art-aware crop come from the rocket source.
The generic crop additionally misclassifies the cat and coffee, each in 30
conditions. Thus crop-based recall gains are not a free accuracy improvement.
For the original alone, generic/art-aware full-image margins are
$-0.00601779/-0.00533542$; the corresponding oracle-artwork margins are
$+0.02146989/+0.02424490$. This diagnostic does not form part of the primary
Luna method, which receives the original image directly.

On the same 18 inputs as Table~\ref{tab:matched}, the generic and art-aware
oracle conditions have TP/FN $13/0$ in both cases, FP/TN $3/2$ and $1/4$,
and accuracies $83.33\%$ and $94.44\%$. Their FPRs are $60\%$ and $20\%$.
The two full-image CLIP FPRs are $40\%$ and $0\%$; their balanced accuracies
are $53.08\%$ and $88.46\%$. The Luna review has BA $100\%$ on this
exploratory set. These figures share a content label, unlike native NSFW scores.

NudeNet's core-body-part proxy detects 0/30 artwork conditions and alarms on
0/150 benign conditions. This is a proxy, not a task-matched nudity accuracy
estimate. The maximum Falconsai and AdamCodd NSFW scores among the 30 artwork
conditions are approximately 0.00294859 and 0.05213738.

\subsection{Visibility measurements}

\begin{table*}[tp]
\centering
\small
\begin{tabular}{@{}lrrr@{}}
\toprule
Scale & Reference area / frame & Falconsai effective pixels & AdamCodd effective pixels \\
\midrule
0.50 & 0.9025\% & 447 & 1334 \\
0.75 & 2.0247\% & 1008 & 2993 \\
1.00 & 3.6076\% & 1816 & 5319 \\
\bottomrule
\end{tabular}
\caption{Coarse reference visibility, not sexual severity. At a fixed scale,
reference-region RGB hashes after each ViT's actual preprocessing agree across
backgrounds. The mask was provided by the research agent.}
\end{table*}

\section{Input identity, runtime, and evidence inventory}
\label{app:repro}

\paragraph{Original-image identity.}
The original JPEG, the figure asset, and Luna input \texttt{S04} all have the
following SHA-256 (shown on two lines):
\begin{quote}\ttfamily\small
93a42b0c2b069cbd9363cf116b0d990d\\
055975376e9cde8951f8a12bcce506ac
\end{quote}

\paragraph{Model provenance.}
The following repositories and revisions were frozen:
\begin{itemize}\small
\item Falconsai export: \path{onnx-community/nsfw_image_detection-ONNX}\\
\texttt{1ceb3c7fe1e9f3f2507e}\allowbreak\texttt{6df577437f23a9149fd5}.
\item AdamCodd: \path{AdamCodd/vit-base-nsfw-detector}\\
\texttt{8587de998f441aac03fd}\allowbreak\texttt{d57a85d2e4cb808c7d64}.
\item CLIP revision: \texttt{d15189d7028b43f1d3e6}\allowbreak\texttt{5039190477f6af591c2a}.
\end{itemize}
The JSON metadata records repository identifiers and SHA-256 hashes of actual
weights and preprocessing configuration. All local ONNX inference uses CPU,
four intra-operation threads, and one inter-operation thread. The ViT input
sizes are $224\times224$ and $384\times384$; NudeNet uses its native 320-pixel
pipeline. The local runtime includes ONNX Runtime 1.24.4, NumPy 1.26.4,
Pillow 12.2.0, and NudeNet 3.4.2.

\paragraph{Review records.}
The research bundle retains \path{plain_results.json},
\path{reference_labels.json}, and \path{inputs/order.json} under
\path{experiments/luna_audit/}. The requested model name is
\texttt{gpt-5.6-luna}, not an image-service alias. The context viewed all 18
images sequentially. We did not collect independent per-image repetitions,
billing latency, or sampling controls. The English manuscript does not imply
that the Chinese review instructions were rerun in English.
An additional context-reminder condition is archived in the research bundle;
it changed none of the 18 content labels and is omitted from the main comparison.

\paragraph{Raw detector and auxiliary records.}
Files \path{image_falconsai.jsonl}, \path{image_adamcodd.jsonl},
\path{image_nudenet.jsonl}, and \path{strategy_scores.jsonl} are stored
under \path{experiments/results/}, alongside \path{image_manifest.json}
and model metadata. There are 211 records per model. Twelve auxiliary images
without an oracle artwork region have no crop result; exports mark those
fields as not run rather than filling them with fabricated scores.

\paragraph{Historical outputs and service availability.}
Two historical restoration outputs were byte-verified but are not independent
source artworks. Their Falconsai scores are 0.00016741 and 0.00015874, and
their AdamCodd scores are 0.00612020 and 0.00463306. The current experiment
generated no new API image: authenticated attempts returned HTTP 429 with
image quota unavailable. This is an infrastructure limitation, not a safety
refusal or an observed attack success. No service credentials appear in the
manuscript or its evidence tables.

\paragraph{Template and draft status.}
This local, unsubmitted manuscript uses the unmodified official ACL style
and bibliography files from \url{https://github.com/acl-org/acl-style-files}.
The template revision and build instructions are supplied locally.
Publication rights for the user-provided image and author metadata must be
resolved before submission or redistribution.

\end{document}
